%%%%%%%%%%%%%%%%%%%%%%%%%%%%%%%%%%%%%%%%%%%%%%%%%%%%%%%%%%%%%%%%%
% Presentacion 1/3: Introduccion a Python para Programadores de C
% Serie "Python desde C" - TP4 Parte B (UNCuyo - Informatica y Programacion)
% Compilar con: pdflatex -shell-escape 01-fundamentos.tex  (dos pasadas)
%%%%%%%%%%%%%%%%%%%%%%%%%%%%%%%%%%%%%%%%%%%%%%%%%%%%%%%%%%%%%%%%%
\documentclass{beamer}

% Paquetes basicos
\usepackage[utf8]{inputenc}
\usepackage[T1]{fontenc}
\usepackage[spanish]{babel}

% Paquete para el resaltado de codigo de alta calidad
\usepackage{minted}

% Tema visual de la presentacion
\usetheme{Madrid}
\usecolortheme{beaver}
\setbeamertemplate{navigation symbols}{}

% Recapitula el indice, resaltando la seccion actual, al empezar cada seccion
\AtBeginSection[]{
  \begin{frame}{Indice de Contenidos}
    \tableofcontents[currentsection]
  \end{frame}
}

% --- Informacion del Titulo ---
\title[Python desde C -- 1/3]{Introducción a Python para Programadores de C}
\subtitle{1/3 -- Sintaxis, tipos, control de flujo y funciones}
\author{Informática y Programación}
\date{\today}
\institute{UNCuyo -- Facultad de Ingeniería}

% --- Inicio del Documento ---
\begin{document}

% --- Diapositiva del Titulo ---
\begin{frame}
  \titlepage
\end{frame}

% --- Diapositiva del Indice ---
\begin{frame}{Índice de Contenidos}
  \tableofcontents
\end{frame}

\begin{frame}{Por qué cambiamos de lenguaje ahora}
  \begin{block}{No es un lenguaje nuevo desde cero}
    Ya conocés variables, tipos, condicionales, bucles y funciones —
    los mismos \textbf{conceptos} de siempre, pero sin compilar y con menos
    sintaxis. Esta serie va a apoyarse todo el tiempo en lo que ya sabés de
    C para aprender Python por \textbf{contraste}, no por repetición.
  \end{block}
  \begin{alertblock}{Esta es la presentación 1 de 3}
    \begin{enumerate}
      \item \textbf{Fundamentos} (ésta): sintaxis, tipos, control, E/S, funciones.
      \item \textbf{Estructuras de datos}: listas, tuplas y diccionarios.
      \item \textbf{Cadenas, archivos y funciones avanzadas}.
    \end{enumerate}
  \end{alertblock}
\end{frame}


% ===============================================================
\section{1. Sintaxis Fundamental}
% ===============================================================

\begin{frame}[fragile]{¡Hola, Mundo!}
  \begin{block}{Puntos Clave}
    \begin{itemize}
      \item Python es un lenguaje interpretado, no requiere compilación explícita.
      \item No hay función \texttt{main()} obligatoria.
      \item El punto y coma \texttt{;} es opcional y no se usa.
      \item Los bloques de código se definen por \textbf{indentación}, no por llaves \texttt{\{\}}.
    \end{itemize}
  \end{block}

  \begin{columns}[T]
    \begin{column}{.48\textwidth}
      \centering\textbf{C}
      \begin{minted}[fontsize=\small, frame=lines, linenos]{c}
#include <stdio.h>

int main() {
    printf("Hola, Mundo!\n");
    return 0;
}
      \end{minted}
    \end{column}
    \begin{column}{.48\textwidth}
      \centering\textbf{Python}
      \begin{minted}[fontsize=\small, frame=lines, linenos]{python}
print("Hola, Mundo!")
      \end{minted}
    \end{column}
  \end{columns}
\end{frame}

\begin{frame}[fragile]{Cuidado: la indentación es sintaxis}
  \begin{alertblock}{No es solo prolijidad}
    En C, la indentación es una convención visual — el compilador no la
    mira, solo le importan las llaves. En Python, la indentación
    \textbf{define} qué líneas pertenecen a qué bloque. Mezclar espacios y
    tabs, o indentar mal, es un error de sintaxis real (\texttt{IndentationError}).
  \end{alertblock}
  \begin{minted}[fontsize=\small, frame=lines]{python}
if 5 > 3:
    print("Se ejecuta")   # pertenece al if (4 espacios)
      print("Error!")     # IndentationError: indentado de mas
  \end{minted}
\end{frame}

% ===============================================================
\section{2. Variables y Tipos de Datos}
% ===============================================================

\begin{frame}[fragile]{Variables y Tipos de Datos}
  \begin{block}{Diferencia Principal: Tipado}
    \begin{itemize}
        \item \textbf{C es de tipado estático:} Debes declarar el tipo de la variable antes de usarla (\texttt{int}, \texttt{float}, etc.).
        \item \textbf{Python es de tipado dinámico:} El tipo se infiere automáticamente en el momento de la asignación.
    \end{itemize}
  \end{block}

  \begin{columns}[T]
    \begin{column}{.48\textwidth}
      \centering\textbf{C}
      \begin{minted}[fontsize=\small, frame=lines, linenos]{c}
int cantidad = 10;
float precio = 95.50;
      \end{minted}
    \end{column}
    \begin{column}{.48\textwidth}
      \centering\textbf{Python}
      \begin{minted}[fontsize=\small, frame=lines, linenos]{python}
# La variable toma el tipo
# del valor asignado
cantidad = 10
precio = 95.50

# Podemos verificar el tipo
# type(cantidad) -> <class 'int'>
      \end{minted}
    \end{column}
  \end{columns}
\end{frame}

\begin{frame}[fragile]{Cuidado: dinámico no es ``sin tipo''}
  \begin{alertblock}{La variable puede cambiar de tipo, el valor no}
    ``Dinámico'' significa que el mismo \textit{nombre} puede referirse a
    valores de distinto tipo en distintos momentos — no que Python ignore
    los tipos. Sigue siendo \textbf{fuertemente tipado}: no suma un
    \texttt{int} y un \texttt{str} sin conversión explícita, igual que C.
  \end{alertblock}
  \begin{minted}[fontsize=\small, frame=lines]{python}
x = 10        # x es int
x = "diez"    # ahora x es str (valido, x cambio de tipo)

edad = 20
print("Tengo " + edad)   # TypeError: no se puede sumar str + int
print("Tengo " + str(edad))  # correcto: conversion explicita
  \end{minted}
\end{frame}

% ===============================================================
\section{3. Estructuras de Control}
% ===============================================================

\begin{frame}[fragile]{Condicionales: if / elif / else}
  \begin{alertblock}{Lógica de Ejecución}
    La estructura ejecuta \textbf{únicamente el primer bloque} cuya condición sea verdadera (\texttt{true}) y luego sale por completo, ignorando el resto.
  \end{alertblock}

  \begin{columns}[T]
    \begin{column}{.48\textwidth}
      \centering\textbf{C}
      \begin{minted}[fontsize=\small, frame=lines, linenos]{c}
int nota = 75;

if (nota >= 90) {
    printf("Excelente\n");
} else if (nota >= 60) {
    printf("Aprobado\n");
} else {
    printf("Reprobado\n");
}
// Salida: Aprobado
      \end{minted}
    \end{column}
    \begin{column}{.48\textwidth}
      \centering\textbf{Python}
      \begin{minted}[fontsize=\small, frame=lines, linenos]{python}
nota = 75

if nota >= 90:
    print("Excelente")
elif nota >= 60:
    print("Aprobado")
else:
    print("Reprobado")
# Salida: Aprobado
      \end{minted}
    \end{column}
  \end{columns}
\end{frame}

\begin{frame}[fragile]{Equivalente a Switch}
  \begin{block}{Simulando un Switch}
      Python no tiene una estructura \texttt{switch-case} tradicional. El comportamiento idéntico se logra con una cadena \texttt{if/elif/else}.
  \end{block}

  \begin{columns}[T]
    \begin{column}{.48\textwidth}
      \centering\textbf{C}
      \begin{minted}[fontsize=\small, frame=lines, linenos]{c}
int dia = 2;
switch (dia) {
    case 1: printf("Lunes\n"); break;
    case 2: printf("Martes\n"); break;
    default: printf("Otro dia\n");
}
      \end{minted}
    \end{column}
    \begin{column}{.48\textwidth}
      \centering\textbf{Python}
      \begin{minted}[fontsize=\small, frame=lines, linenos]{python}
dia = 2
if dia == 1:
    print("Lunes")
elif dia == 2:
    print("Martes")
else:
    print("Otro dia")
      \end{minted}
    \end{column}
  \end{columns}
\end{frame}

\begin{frame}[fragile]{Bucles: while}
  \begin{block}{Sintaxis Similar}
    La lógica del \texttt{while} es idéntica. Cambian los detalles: sin paréntesis en la condición y sin el operador \texttt{i++}.
  \end{block}
  \begin{columns}[T]
    \begin{column}{.48\textwidth}
      \centering\textbf{C}
      \begin{minted}[fontsize=\small, frame=lines, linenos]{c}
int i = 0;
while (i < 5) {
    printf("%d\n", i);
    i++;
}
      \end{minted}
    \end{column}
    \begin{column}{.48\textwidth}
      \centering\textbf{Python}
      \begin{minted}[fontsize=\small, frame=lines, linenos]{python}
i = 0
while i < 5:
    print(i)
    i += 1
      \end{minted}
    \end{column}
  \end{columns}
  \vspace{0.3em}
  \small \textit{Python no tiene \texttt{i++} ni \texttt{i- -}: se usa \texttt{i += 1}.}
\end{frame}


\begin{frame}[fragile]{Bucles: for}
  \begin{block}{Filosofía Diferente}
    El \texttt{for} de C es un contador. El \texttt{for} de Python itera sobre secuencias. Para replicar el comportamiento de C, usamos \texttt{range()}.
  \end{block}
  \begin{columns}[T]
    \begin{column}{.48\textwidth}
      \centering\textbf{C (Contador)}
      \begin{minted}[fontsize=\small, frame=lines, linenos]{c}
int i;
for (i = 0; i < 5; i++) {
    printf("%d\n", i);
}
      \end{minted}
    \end{column}
    \begin{column}{.48\textwidth}
      \centering\textbf{Python (Iterador)}
      \begin{minted}[fontsize=\small, frame=lines, linenos]{python}
# range(5) genera 0, 1, 2, 3, 4
for i in range(5):
    print(i)
      \end{minted}
    \end{column}
  \end{columns}
\end{frame}

\begin{frame}[fragile]{\texttt{range()} con inicio, fin y paso}
  \begin{block}{Igual de flexible que un \texttt{for} de C}
    \texttt{range(inicio, fin, paso)} — \texttt{fin} nunca se incluye, igual que la condición de corte de un \texttt{for} en C.
  \end{block}
  \begin{columns}[T]
    \begin{column}{.48\textwidth}
      \centering\textbf{C}
      \begin{minted}[fontsize=\small, frame=lines, linenos]{c}
for (int i = 10; i > 0; i--) {
    printf("%d\n", i);
}
      \end{minted}
    \end{column}
    \begin{column}{.48\textwidth}
      \centering\textbf{Python}
      \begin{minted}[fontsize=\small, frame=lines, linenos]{python}
for i in range(10, 0, -1):
    print(i)
      \end{minted}
    \end{column}
  \end{columns}
\end{frame}

% ===============================================================
\section{4. Entrada de Datos desde Teclado}
% ===============================================================

\begin{frame}[fragile]{Lectura de Datos: scanf vs input()}
  \begin{block}{Diferencia Clave}
    \begin{itemize}
      \item En C usas \texttt{scanf()} con especificadores de formato.
      \item En Python usas \texttt{input()}, que \textbf{siempre retorna un string}.
      \item Debes convertir explícitamente a otros tipos (\texttt{int()}, \texttt{float()}).
    \end{itemize}
  \end{block}

  \begin{columns}[T]
    \begin{column}{.48\textwidth}
      \centering\textbf{C}
      \begin{minted}[fontsize=\footnotesize, frame=lines, linenos]{c}
char nombre[50];
int edad;

printf("Nombre: ");
scanf("%s", nombre);

printf("Edad: ");
scanf("%d", &edad);

printf("Hola %s\n", nombre);
      \end{minted}
    \end{column}
    \begin{column}{.48\textwidth}
      \centering\textbf{Python}
      \begin{minted}[fontsize=\footnotesize, frame=lines, linenos]{python}
# input() retorna string
nombre = input("Nombre: ")

# Convertimos a entero
edad = int(input("Edad: "))

print(f"Hola {nombre}")
      \end{minted}
    \end{column}
  \end{columns}
\end{frame}

\begin{frame}[fragile]{Lectura de Múltiples Valores}
  \begin{block}{Técnica con split()}
    Python puede leer varios valores en una línea usando \texttt{split()} y desempaquetado.
  \end{block}

  \begin{columns}[T]
    \begin{column}{.48\textwidth}
      \centering\textbf{C}
      \begin{minted}[fontsize=\footnotesize, frame=lines, linenos]{c}
int a, b;

printf("Dos numeros: ");
scanf("%d %d", &a, &b);

printf("Suma: %d\n", a + b);
      \end{minted}
    \end{column}
    \begin{column}{.48\textwidth}
      \centering\textbf{Python}
      \begin{minted}[fontsize=\footnotesize, frame=lines, linenos]{python}
# Leer y separar
entrada = input("Dos numeros: ")
a, b = entrada.split()

# Convertir a enteros
a = int(a)
b = int(b)

print(f"Suma: {a + b}")
      \end{minted}
    \end{column}
  \end{columns}
\end{frame}

% ===============================================================
\section{5. Funciones y Parámetros}
% ===============================================================

\begin{frame}[fragile]{Definición de Funciones}
  \begin{block}{Definición Simplificada}
    En Python se usa la palabra clave \texttt{def}. No se declaran tipos de retorno ni de parámetros.
  \end{block}
  \begin{columns}[T]
    \begin{column}{.48\textwidth}
      \centering\textbf{C}
      \begin{minted}[fontsize=\small, frame=lines, linenos]{c}
int sumar(int a, int b) {
    return a + b;
}

int main() {
    int res = sumar(10, 5);
    // res es 15
    return 0;
}
      \end{minted}
    \end{column}
    \begin{column}{.48\textwidth}
      \centering\textbf{Python}
      \begin{minted}[fontsize=\small, frame=lines, linenos]{python}
def sumar(a, b):
    return a + b

# --- Llamada ---
res = sumar(10, 5)
# res es 15
      \end{minted}
    \end{column}
  \end{columns}
\end{frame}

\begin{frame}{Paso de Parámetros: Mutabilidad}
  \begin{alertblock}{Punto Clave}
    \begin{itemize}
      \item \textbf{C pasa copias (por valor).} Para modificar el original, se usan punteros (\texttt{*}).
      \item \textbf{Python pasa referencias a objetos.}
      \begin{itemize}
        \item Objetos \textbf{inmutables} (números, strings) se comportan como paso por valor.
        \item Objetos \textbf{mutables} (listas, diccionarios) se comportan como paso por referencia.
      \end{itemize}
      \item El resultado final en el ejemplo que sigue es idéntico en los dos lenguajes.
    \end{itemize}
  \end{alertblock}
\end{frame}

\begin{frame}[fragile]{Paso de Parámetros: mismo resultado en los dos lenguajes}
  \begin{columns}[T]
    \begin{column}{.48\textwidth}
      \centering\textbf{C}
      \begin{minted}[fontsize=\footnotesize, frame=lines, linenos]{c}
void mod(int num, int arr[]) {
    num = 100;   // no cambia
    arr[0] = 99; // SI cambia
}

// mi_numero -> 10
// mi_array[0] -> 99
      \end{minted}
    \end{column}
    \begin{column}{.48\textwidth}
      \centering\textbf{Python}
      \begin{minted}[fontsize=\footnotesize, frame=lines, linenos]{python}
def mod(num, lista):
    num = 100     # no cambia
    lista[0] = 99 # SI cambia

# mi_numero -> 10
# mi_lista[0] -> 99
      \end{minted}
    \end{column}
  \end{columns}
  \vspace{0.4em}
  {\footnotesize \textit{Por qué se parecen tanto: en C, un arreglo pasado
  como parámetro ya decae a puntero — la \emph{dirección} se copia por
  valor. Reasignar \texttt{arr} adentro de la función no movería el
  arreglo del \texttt{main}, igual que reasignar \texttt{lista} en Python
  no afecta a la variable de afuera. Lo que sí se comparte, en los dos
  lenguajes, es lo que hay \emph{en} esa dirección.}}
\end{frame}

% ===============================================================
\section{6. Resumen y Conclusiones}
% ===============================================================

\begin{frame}{Tabla Comparativa Rápida}
  \begin{tabular}{|l|p{4cm}|p{4cm}|}
    \hline
    \textbf{Característica} & \textbf{Lenguaje C} & \textbf{Lenguaje Python} \\
    \hline \hline
    Tipado & Estático (fuerte) & Dinámico (fuerte) \\
    \hline
    Compilación & Compilado & Interpretado \\
    \hline
    Sintaxis & Verbosa, con \texttt{;} y \texttt{\{\}} & Concisa, basada en indentación \\
    \hline
    Manejo de Memoria & Manual (\texttt{malloc}, \texttt{free}) & Automático (Garbage Collector) \\
    \hline
    Punteros & Explícitos y fundamentales & No existen de forma directa \\
    \hline
    Bucle \texttt{for} & Basado en un contador & Basado en iteración sobre secuencias \\
    \hline
  \end{tabular}
\end{frame}

\begin{frame}{Lo que sigue}
  \begin{block}{Con esto ya podés}
    Traducir a Python cualquier programa de C que use solo variables,
    condicionales, bucles y funciones con parámetros simples — la base de
    los ejercicios 5 y 6 del TP4 (parking, batalla naval).
  \end{block}
  \begin{alertblock}{Presentación 2: Estructuras de Datos}
    Para modelar los tres pisos del parking, el tablero de la batalla
    naval, o el inventario de la librería, hace falta más que variables
    sueltas: \textbf{listas, tuplas y diccionarios} — el equivalente
    Python de los arreglos y \texttt{structs} de C. Eso es lo que viene.
  \end{alertblock}
\end{frame}

\begin{frame}
  \vfill
  \centering
  \begin{beamercolorbox}[sep=8pt,center,shadow=true,rounded=true]{title}
    \Huge \bfseries ¡A programar en Python!
  \end{beamercolorbox}
  \vfill
\end{frame}

\end{document}
